% Content for the Innovation in Academia brochure.
\colorlet{Accent}{Coral}
\WorkshopHeader
  {01}
  {INNOVATION IN ACADEMIA}
  {A faculty transformation lab for stronger patents, funded research and industry relevance}
  {Research novelty \textbullet\ Patent potential \textbullet\ Fundable partnerships}
  {\faLightbulb}
  {assets/bg-innovation-academia.png}

\MetaStrip
  {Teachers \& academic leaders}
  {1 day \textbar\ 9:30--17:00}
  {Hands-on studio}
  {INR 3,500 / person}

\begin{minipage}[t]{0.615\textwidth}
  \vspace{0pt}
  \begin{brochurecard}[colback=Accent!3!white]
    \SectionHead{Why this matters now}
    \LeaderHook{Turn faculty expertise into a visible pipeline of novel research, defensible IP and fundable partnerships.}
    \WhyItMatters
      {Research scholars learn to identify consequential problems, build inventive confidence and see how research becomes IP and real-world impact.}
      {Faculty move beyond incremental publication, strengthen novelty and create credible paths to patents, funded research and industry partnerships.}
  \end{brochurecard}

  \begin{brochurecard}
    \begin{minipage}[t]{0.485\linewidth}
      \SectionHead{What we will explore}
      \begin{itemize}
        \item Build an inventive mindset for research, teaching and institutional challenges.
        \item Separate genuine novelty from routine extensions of existing work.
        \item Use structured frameworks to connect evidence, IP and stakeholder value.
      \end{itemize}
    \end{minipage}\hfill
    \begin{minipage}[t]{0.485\linewidth}
      \SectionHead{What participants gain}
      \begin{itemize}
        \item \textbf{Find novelty:} identify innovative, relevant problems worth solving.
        \item \textbf{Escape the publication trap:} move from incremental papers to defensible contribution.
        \item \textbf{Create IP:} shape novel, non-obvious and patent-ready concepts.
        \item \textbf{Attract support:} build credible industry and government-funding pathways.
      \end{itemize}
    \end{minipage}
  \end{brochurecard}

  \begin{brochurecard}
    \SectionHead{One-day learning journey}
    \ScheduleRow{09:30--10:15}{\textbf{Beyond the next paper:} recognise the iterative publication trap and define meaningful impact.}
    \ScheduleRow{10:15--11:15}{\textbf{Novel problem discovery:} scan tensions, unmet needs, anomalies and emerging opportunity landscapes.}
    \ScheduleRow{11:30--12:45}{\textbf{Storytelling for relevance:} connect evidence, beneficiaries, urgency and a compelling research question.}
    \ScheduleRow{13:45--15:00}{\textbf{Framework lab:} TRIZ contradictions, inventive principles and structured idea generation.}
    \ScheduleRow{15:15--16:15}{\textbf{Reverse brainstorming and IP lens:} expose failure modes and test novelty and non-obviousness.}
    \ScheduleRow{16:15--17:00}{\textbf{Impact pathway:} map experiments, patent steps, industry partners and suitable funding routes.}
  \end{brochurecard}
\end{minipage}\hfill
\begin{minipage}[t]{0.36\textwidth}
  \vspace{0pt}
  \begin{brochurecard}[colback=Accent!5!white]
    \SectionHead{Who should attend}
    \begin{itemize}
      \item University and college teachers
      \item Heads of department and deans
      \item Research supervisors and lab leaders
      \item Curriculum, accreditation and faculty-development teams
    \end{itemize}
  \end{brochurecard}

  \begin{brochurecard}
    \SectionHead{Methods \& tools}
    \begin{itemize}
      \item Novel problem and opportunity landscapes
      \item Research storytelling for relevance
      \item TRIZ contradictions and principles
      \item Reverse brainstorming
      \item Publication-to-patent decision lens
      \item Industry collaboration and funding canvas
    \end{itemize}
  \end{brochurecard}

  \begin{brochurecard}[colback=Accent!3!white]
    \SectionHead{Takeaway toolkit}
    \begin{itemize}
      \item A mapped portfolio of consequential, underexplored research problems
      \item A novelty and non-obviousness canvas for a selected idea
      \item Concepts generated using TRIZ and reverse brainstorming
      \item A 30-day experiment, IP and stakeholder-engagement roadmap
    \end{itemize}
    {\color{Accent}\bfseries\fontsize{6.6}{7.2}\selectfont INSTITUTIONAL RETURN}\par
    A stronger pipeline of patentable research, fundable proposals and industry-facing faculty initiatives.
  \end{brochurecard}

  \begin{feecard}
    \SectionHead{Investment \& included value}
    {\bfseries INR 3,500 per participant}\par
    Workbook, innovation tool cards, IP lens, funding canvas and participation certificate included.
    Institutional cohort pricing available for 25--40 participants.

    {\color{Slate}\fontsize{6.8}{8}\selectfont
      Indicative fee; taxes, travel and venue costs are additional where applicable.}
  \end{feecard}
\end{minipage}

\Footer